\chapter{Quantifier Elimination and Categoricity}
\label{ch:qe-categoricity}

\section{Quantifier elimination}
\label{sec:qe-categoricity-01}

\begin{definition}[Quantifier elimination]
A theory \(T\) in a signature \(\Lang\) \emph{eliminates quantifiers} iff,
for every \(\Lang\)-formula \(\varphi(\bar x)\), there is a quantifier-free
\(\Lang\)-formula \(\psi(\bar x)\) such that
\[
  T\entails\forall\bar x\,
  \bigl(\varphi(\bar x)\leftrightarrow\psi(\bar x)\bigr).
\]
The equivalence is relative to the models of \(T\); it is not asserted in all
\(\Lang\)-structures.
\end{definition}

\begin{method}[One-quantifier reduction]
To prove quantifier elimination, it suffices to give an explicit reduction of
every formula
\(\exists y\,\theta(y,\bar x)\), with \(\theta\) quantifier-free, to a
quantifier-free formula modulo \(T\).  Boolean induction, replacement of
\(\forall y\,\theta\) by \(\neg\exists y\,\neg\theta\), and renaming of
bound variables then remove all quantifiers from arbitrary formulas.
\end{method}

\begin{theorem}[Dense linear orders]
\label{thm:dlo-qe}
Let \(\mathsf{DLO}\) be the first-order theory of nonempty dense linear orders
without endpoints in the signature \(\{<\}\).  Then \(\mathsf{DLO}\)
eliminates quantifiers.
\textit{Status: cited theorem.}
\end{theorem}

\begin{remark}[Construction underlying the cited theorem]
Put a quantifier-free matrix in disjunctive normal form and distribute the
existential quantifier over the finite disjunction.  In one conjunction,
separate literals involving \(y\) into lower bounds \(t_i<y\), upper bounds
\(y<r_j\), equalities \(y=u_k\), and disequalities.  An equality substitutes
its term for \(y\).  Without an equality, density and the absence of endpoints
show that a witness exists iff every required lower bound is below
every required upper bound, subject to the finitely many equalities and
disequalities among parameters.  Writing those finite compatibility conditions
gives a quantifier-free formula.  A complete proof verifies all inconsistent
literal combinations and then applies the one-quantifier reduction method; see
\cite{Marker2002}.
\end{remark}

\begin{example}[Eliminating one quantifier in \(\mathsf{DLO}\)]
Modulo \(\mathsf{DLO}\),
\[
  \exists y\,(x<y\land y<z)
  \quad\longleftrightarrow\quad x<z.
\]
The forward implication uses transitivity.  The reverse implication is exactly
density.
\end{example}

\begin{definition}[Deductive completeness, semantic completeness, and model completeness]
\label{def:theory-completeness}
A theory \(T\) is \emph{deductively complete} iff, for every sentence
\(\sigma\) in its signature, either \(T\derives\sigma\) or
\(T\derives\neg\sigma\).  It is \emph{semantically complete} iff, for every
such \(\sigma\), either \(T\entails\sigma\) or
\(T\entails\neg\sigma\).  Because derivability in this reference is taken
in the
first-order calculus fixed by
\cref{thm:fol-soundness,thm:fol-completeness},
\[
  T\derives\sigma\quad\Longleftrightarrow\quad T\entails\sigma,
\]
so these two properties are equivalent here.  The fixed term for the
derivability property is ``deductively complete,'' and the fixed term for the
semantic property is ``semantically complete.''

A theory is \emph{model complete} iff every embedding between its models is
elementary.  This is not either decision property above: it concerns formulas
with parameters along embeddings, and later results require that distinction.
\end{definition}

\begin{lemma}[Reduction to substructure inclusions]
\label{lem:model-completeness-reduction}
Let \(T\) be a theory such that
\(\mathcal M\preccurlyeq\mathcal N\) whenever
\(\mathcal M\subseteq\mathcal N\) are both models of \(T\).  Then \(T\)
is model complete.
\textit{Status: complete proof.}
\end{lemma}
\begin{proof}
Let \(h:\mathcal M\to\mathcal N\) be an embedding between models of
\(T\).  By \cref{lem:embedding-image}, \(h\) factors as the isomorphism
\(h_0:\mathcal M\to h[\mathcal M]\) obtained by corestriction, followed
by the inclusion \(h[\mathcal M]\subseteq\mathcal N\).  By
\cref{lem:iso-invariance}, \(h[\mathcal M]\) satisfies the same sentences
as \(\mathcal M\), so \(h[\mathcal M]\) is a model of \(T\).  The
hypothesis applied to \(h[\mathcal M]\subseteq\mathcal N\) gives
\(h[\mathcal M]\preccurlyeq\mathcal N\).  Then, for every formula
\(\varphi(x_1,\ldots,x_k)\) and every tuple \(\bar a\) from \(M\),
\[
  \mathcal M\entails\varphi(\bar a)
  \quad\Longleftrightarrow\quad
  h[\mathcal M]\entails\varphi(h\bar a)
  \quad\Longleftrightarrow\quad
  \mathcal N\entails\varphi(h\bar a),
\]
the first equivalence by \cref{lem:iso-invariance} applied to \(h_0\) and
the second because the inclusion is elementary.  Hence \(h\) is
elementary, and \(T\) is model complete.
\end{proof}

\begin{corollary}
The theory \(\mathsf{DLO}\) is deductively complete (equivalently,
semantically complete) and model complete.
\textit{Status: proof sketch.}
\end{corollary}

\begin{proofplan}
A sentence is equivalent modulo \(\mathsf{DLO}\) to a quantifier-free sentence.
The signature has no constant symbols, so there are no closed atomic
formulas; every quantifier-free sentence is a truth-functional combination
of \(\bottom\), and its truth value is fixed by the semantics of the
connectives alone.  Hence any two models of
\(\mathsf{DLO}\) satisfy the same sentences.  Thus the theory is semantically
complete and therefore deductively complete by
\cref{def:theory-completeness}.
If \(\mathcal M\subseteq\mathcal N\) are models of \(\mathsf{DLO}\), they agree
on quantifier-free formulas with parameters from \(M\).  Quantifier elimination
therefore gives agreement on all formulas, so
\(\mathcal M\preccurlyeq\mathcal N\).  By
\cref{lem:model-completeness-reduction}, every embedding between models
of \(\mathsf{DLO}\) is elementary; this is model completeness.
The omitted literal-case verification is the one described in the remark
following \cref{thm:dlo-qe}.
\end{proofplan}

\begin{theorem}[Algebraically closed fields]
For each characteristic \(p\), where \(p=0\) or \(p\) is prime, the theory
\(\mathsf{ACF}_p\) of algebraically closed fields of characteristic \(p\)
eliminates quantifiers in the language of rings.
\textit{Status: cited theorem.}
\end{theorem}

See \cite{Marker2002} for a complete proof.

